\documentclass[10pt,a4paper]{amsart}
\usepackage[margin=19mm]{geometry}
\usepackage{amsmath,amssymb,mathtools}
\usepackage[scr=rsfs,cal=euler]{mathalfa}
\usepackage{xfrac}
\usepackage[unicode,hidelinks,bookmarks=false]{hyperref}
\newcommand{\ie}{i.e.\ }
\newcommand{\cf}{cf.\ }
\newcommand{\resp}{respectively\ }
\newcommand{\Subsection}{\subsection}
\providecommand{\nobreakdash}{\nobreak\hbox{-}}
\setlength{\emergencystretch}{2em}
\multlinegap0pt
\usepackage{bm}
\usepackage{bbm}
\usepackage{dsfont}
\usepackage{xspace}
\usepackage{cancel}
\usepackage{mathtools}
\usepackage{cleveref}
\usepackage{amsthm}
\makeatletter
\@ifundefined{theorem}{\newtheorem{theorem}{Theorem}}{}
\@ifundefined{lemma}{\newtheorem{lemma}[theorem]{Lemma}}{}
\@ifundefined{proposition}{\newtheorem{proposition}[theorem]{Proposition}}{}
\makeatother
\newcommand{\E}{\mathbb{E}}
\newcommand{\R}{\mathbb{R}}
\newcommand{\eps}{\varepsilon}
\newcommand{\mb}{\mathbb}
\newcommand{\mc}{\mathcal}
\renewcommand{\P}{\mathbb{P}}
\title[The online monotone array completion problem]{The online monotone array completion problem}
\author{Vishesh Jain, Dylan King, Clayton Mizgerd}
\date{Source: arXiv:2606.32015v1 (CC BY 4.0)}
\begin{document}
\maketitle

\noindent\emph{Source and licence.} The online monotone array completion problem. Version arXiv:2606.32015v1 (CC BY 4.0), distributed under the Creative Commons Attribution 4.0 International licence, as recorded in the arXiv licence metadata for this exact version and shown on the abstract page https://arxiv.org/abs/2606.32015v1. Full rights notice: licenses/arxiv-2606.32015-CC-BY-4.0.txt.\par\smallskip

\noindent\textit{Editorial scope.} Counted unit: the Proposition whose source label is \texttt{prop:upper-explicit} in the article (source0.tex lines 1175--1185, source label \texttt{prop:upper-explicit}), delivered together with the article's own complete original proof of it (source0.tex lines 1187--1342, one \texttt{proof} environment of 156 lines). No mathematics is written, repaired or re-derived: statement and proof are the article's own text line for line. Required context retained uncounted: lem:feasible-length (lines 1109--1136). Every definition, display and label of those lines is retained unchanged. Omitted from this package and not redistributed: 1--1108, 1137--1174, 1186--1186, 1343--1780. Nothing inside a retained excerpt is omitted or altered: every retained body is delivered exactly as the article's lines are written, so no omission receipt of any kind accompanies this package. Citations: the retained excerpts carry no citation command at all, so no bibliography entry of the article is transcribed and no reference is redistributed. Reference adaptation: none was needed and none was made; every reference inside the delivered unit resolves inside it, so no unresolved reference or citation is produced. Printed slips kept literal: no printed word, symbol, hypothesis or number of the counted statement or of its proof was altered. Layout and class: the article's own class is \texttt{amsart} with packages including booktabs, color, hyperref, amsmath, amssymb, amsthm, mathrsfs, mathtools, cleveref, fullpage, cite, graphics, pifont, tikz; the wrapper is the lane's pinned \texttt{amsart} editorial wrapper at 10pt on A4, which restates the theorem-like environments the retained text needs and adds the article's own macro definitions that the retained text needs (\texttt{\textbackslash E}, \texttt{\textbackslash P}, \texttt{\textbackslash R}, \texttt{\textbackslash eps}, \texttt{\textbackslash mb}, \texttt{\textbackslash mc}), with extra wrapper packages (bm, bbm, dsfont, xspace, cancel, mathtools, cleveref). The delivered numbering is the unit's own. Assets: no retained span contains a figure environment, a graphics-inclusion command, a PDF-page inclusion, a table or a TikZ picture; no image file is copied into this package, the package declares no asset dependency and no image file is redistributed with it. Licence: version arXiv:2606.32015v1 of this article is distributed under the Creative Commons Attribution 4.0 International licence, as recorded in the arXiv licence metadata for this exact version and shown on the abstract page https://arxiv.org/abs/2606.32015v1. A full rights notice is delivered at licenses/arxiv-2606.32015-CC-BY-4.0.txt. Characters: every retained excerpt is pure ASCII, so the delivered engine, XeLaTeX with its default Latin Modern text font, reports no missing character.
\begin{lemma}\label{lem:feasible-length}
Fix a block \(B_j\) of size \(s=s_j\). Suppose that \(k\ge1\) locations remain unfilled in this block, and write the current feasible interval as
\[
I_j^{\rm feas}=[L,R],
\qquad
\ell:=R-L.
\]
Then
\[
\ell\ge \frac{k+1}{s+1}\cdot \frac{s}{n}.
\]
Consequently, if
\[
E^-:=\left[L,L+\frac{\ell}{k+1}\right],
\qquad
E^+:=\left[R-\frac{\ell}{k+1},R\right]
\]
are the two edge intervals from which the strategy accepts samples, then
\[
|E^-|+|E^+|
=
\frac{2\ell}{k+1}
\ge
\frac{2s}{(s+1)n}
\ge
\frac{2b}{(b+1)n}.
\]
\end{lemma}

\begin{proposition}\label{prop:upper-explicit}
For each \(n\), there is a deterministic strategy \(S=S^{(n)}\) such that
\[
\mb{E}[\tau_{S}]\le \left(\frac12+o(1)\right)n\log n.
\]
Moreover, for every fixed \(\eps>0\),
\[
\P\left(\tau_{S}>\left(\frac12+\eps\right)n\log n\right)\le n^{-\eps}
\]
for all sufficiently large \(n\).
\end{proposition}

\begin{proof}
Set
\[
b:=\left\lceil \sqrt{\log n}\right\rceil
\]
and define \(S^{(n)}:=S_b\). Let \(\tau_b\) denote the completion time of this strategy, and set
\[
\lambda:=\frac{2b}{(b+1)n}.
\]
By \cref{lem:feasible-length}, whenever a block is unfinished, the current acceptance region for that block has length at least \(\lambda\).

Fix a block \(B_j\), and let \(\tau_j\) be the time at which \(B_j\) is completed. If \(s_j=|B_j|\), write
\[
\tau_j=W_1^{(j)}+\cdots+W_{s_j}^{(j)},
\]
where \(W_r^{(j)}\) is the number of samples between the \((r-1)\)-st and \(r\)-th accepted samples in \(B_j\). Condition on the history up to the time at which \(r-1\) samples have been accepted in \(B_j\). Until the next acceptance in \(B_j\), every new sample has conditional probability at least \(\lambda\) of being accepted by this block. Therefore \(W_r^{(j)}\) is conditionally dominated by a geometric random variable \(G_\lambda\) with
\[
\P(G_\lambda=q)=(1-\lambda)^{q-1}\lambda,
\qquad q\ge1.
\]
Consequently, for every \(0<\theta<-\log(1-\lambda)\),
\[
\E\left[e^{\theta W_r^{(j)}}\mid \mc F_r \right]
\le
\E e^{\theta G_\lambda}
=
\frac{\lambda e^\theta}{1-(1-\lambda)e^\theta},
\]
where~$\mc F_r$ is the natural filtration on all past events.
Iterating over \(r=1,\dots,s_j\), and using \(s_j\le b+1\), gives
\begin{equation}\label{eq:block-mgf-discrete}
\E e^{\theta\tau_j}
\le
\left(\frac{\lambda e^\theta}{1-(1-\lambda)e^\theta}\right)^{b+1}.
\end{equation}

Now choose
\[
\eta:=\frac1{\log n},
\qquad
\theta:=-\log\bigl(1-(1-\eta)\lambda\bigr).
\]
Then
\[
e^\theta=\frac{1}{1-(1-\eta)\lambda},
\]
and hence
\[
\frac{\lambda e^\theta}{1-(1-\lambda)e^\theta}=\frac1\eta.
\]
Thus \eqref{eq:block-mgf-discrete} gives
\begin{equation}\label{eq:block-mgf-specialized}
\E e^{\theta\tau_j}\le \eta^{-(b+1)}
\end{equation}
for every block \(B_j\).

Since the array is complete exactly when all blocks are complete,
\[
\tau_b=\max_{1\le j\le m}\tau_j.
\]
Using
\[
\max_{1\le j\le m} x_j
\le
\frac1\theta\log\left(\sum_{j=1}^m e^{\theta x_j}\right),
\qquad x_1,\dots,x_m\in\R,
\]
together with Jensen's inequality and \eqref{eq:block-mgf-specialized}, we obtain
\[
\E[\tau_b]
\le
\frac1\theta
\log\left(\sum_{j=1}^m \E e^{\theta\tau_j}\right)
\le
\frac1\theta
\left(
\log m+(b+1)\log\frac1\eta
\right).
\]
Since
\[
\theta
=
-\log\bigl(1-(1-\eta)\lambda\bigr)
\ge
(1-\eta)\lambda
=
(1-\eta)\frac{2b}{(b+1)n},
\]
it follows that
\[
\E[\tau_b]
\le
\frac{(b+1)n}{2b(1-\eta)}
\left(
\log m+(b+1)\log\frac1\eta
\right).
\]
For our choices \(b=\lceil\sqrt{\log n}\rceil\) and \(\eta=1/\log n\), we have
\[
\frac{b+1}{b}=1+o(1),
\qquad
\frac1{1-\eta}=1+o(1),
\qquad
\log m=\log n+o(\log n),
\]
and
\[
(b+1)\log\frac1\eta
=
O(\sqrt{\log n}\log\log n)
=
o(\log n).
\]
Therefore
\[
\E[\tau_b]
\le
\left(\frac12+o(1)\right)n\log n.
\]

It remains to prove the high-probability bound. By Markov's inequality and \eqref{eq:block-mgf-specialized}, for every \(t\ge0\),
\begin{equation}\label{eq:block-tail-specialized}
\P(\tau_b>t)
\le
e^{-\theta t}\sum_{j=1}^m \E e^{\theta\tau_j}
\le
\exp\left(
\log m+(b+1)\log\frac1\eta-\theta t
\right).
\end{equation}
Let
\[
t_n:=\left(\frac12+\eps\right)n\log n.
\]
Using \(\theta\ge (1-\eta)\lambda\), \eqref{eq:block-tail-specialized} gives
\begin{align*}
\log \P(\tau_b>t_n)
&\le
\log m+(b+1)\log\frac1\eta
-
(1-\eta)\frac{2b}{(b+1)n}
\left(\frac12+\eps\right)n\log n  \\
&\le
\log n+o(\log n)
-
(1+2\eps)(1-o(1))\log n  \\
&=
(-2\eps+o(1))\log n.
\end{align*}
For all sufficiently large \(n\), this is at most \(-\eps\log n\). Hence
\[
\P\left(\tau_b>\left(\frac12+\eps\right)n\log n\right)\le n^{-\eps},
\]
as claimed.
\end{proof}

\end{document}
